% !PRINTED 305
\weberpage{305}

If we consider two successive terms of this double sequence, such as
\[
\begin{array}{l}
f_\nu,\ f_{\nu+1} \\
F_\nu,\ F_{\nu+1}
\end{array}
\]
for a value of $x$ at which none of the four functions vanishes, then with
regard to the succession of signs the following cases are possible:
\[
\text{a)}\quad
\begin{array}{c}+\ +\\+\ +\end{array},\
\begin{array}{c}+\ +\\-\ -\end{array},\
\begin{array}{c}-\ -\\+\ +\end{array},\
\begin{array}{c}-\ -\\-\ -\end{array},\quad PP,
\]
\[
\text{b)}\quad
\begin{array}{c}+\ -\\+\ -\end{array},\
\begin{array}{c}+\ -\\-\ +\end{array},\
\begin{array}{c}-\ +\\+\ -\end{array},\
\begin{array}{c}-\ +\\-\ +\end{array},\quad VV,
\]
\[
\text{c)}\quad
\begin{array}{c}+\ +\\+\ -\end{array},\
\begin{array}{c}+\ +\\-\ +\end{array},\
\begin{array}{c}-\ -\\+\ -\end{array},\
\begin{array}{c}-\ -\\-\ +\end{array},\quad PV,
\]
\[
\text{d)}\quad
\begin{array}{c}+\ -\\+\ +\end{array},\
\begin{array}{c}+\ -\\-\ -\end{array},\
\begin{array}{c}-\ +\\+\ +\end{array},\
\begin{array}{c}-\ +\\-\ -\end{array},\quad VP.
\]
We denote this behaviour by the following expressions and symbols:

a) succession-succession (permanence-permanence) $PP$,

b) change-change (variation-variation) $VV$,

c) succession-change (permanence-variation) $PV$,

d) change-succession (variation-permanence) $VP$.

The two theorems that we wish to prove read:

\gesperrt{II. The number of roots of $f(x)$ lying between $\alpha$ and
$\beta$ is either exactly as great as, or smaller by an even number than,
the number of change-successions lost in the double sequence (7) in the
passage from $\alpha$ to $\beta$.}

\gesperrt{III. The number of roots of $f(x)$ lying between $\alpha$ and
$\beta$ is either exactly as great as, or smaller by an even number than,
the number of succession-successions gained in the double sequence (7) in
the passage from $\alpha$ to $\beta$.}

Of these two theorems, which do not always give the same upper bound for
the number of roots, one will apply the one that gives the lowest bound.
